\lecture{23}{May 29.}{}
Now since
\[
    \chi_{\ell }^{\prime} (K_{n, n}) \ge \chi^{\prime}  (K_{n, n}) \ge \Delta (K_{n, n}) = n.
\] 
So it suffices to show \(\chi _\ell ^{\prime} (K_{n, n}) \le n\). Therefore, it remains to construct a kernel-perfect orientation \(\vv{D}\) of \(L(K_{n, n})\) with \(\Delta ^+ (\vv{D}) = n - 1\). If so, then 
\[
    \chi _\ell ^{\prime} (K_{n, n}) = \chi _\ell (L(K_{n, n})) \le \Delta ^+(\vv{D}) + 1 = n.
\]  
Hence,
\[
    n \ge \chi _\ell ^{\prime} (K_{n, n}) \ge \chi ^{\prime} (K_{n, n}) \ge n.
\]
This shows \(\chi _\ell ^{\prime} (K_{n, n}) = \chi ^{\prime} (K_{n, n})\).

Now if we can find such kernel-perfect orientation, we can prove the special case of \(K_{n, n}\).

\begin{observation}
    If \(\vv{D}\) is kernel-perfect, then every clique is transitive. (i.e. \(x \to y\) and \(y \to z\) implies \(x \to z\)). 
\end{observation}  

\begin{proof}
    Let \(S\) be the clique. Then \(\vv{D} [S]\) has a kernel. Since \(S\) is a clique and the kernel is independent, the kernel is a single vertex \(u_1\). Threfore, for all \(v \in S \setminus \left\{ u_1 \right\} \), we must have \(v \to u_1\). Repeating for \(S \setminus \left\{ u_1 \right\} \) and so on, we have
    \[
        u_1 \leftarrow u_2 \leftarrow \dots \leftarrow u_{\vert S \vert },
    \]       
    where for any \(i < j\) we have \(u_i \leftarrow u_j\). Thus, if \(u_x \to u_y\) and \(u_y \to u_z\), then we have \(x > y\) and \(y > z\), and thus \(x > z\), which gives \(u_x \to u_z\), which implies the transitivity.        
\end{proof}

\begin{claim}
    There is an orientation \(\vv{D}\) of \(L(K_{n, n})\) such that 
    \begin{enumerate}[(i)]
        \item \(\Delta ^+ (\vv{D}) = n - 1\). 
        \item For every vertex \(v \in K_{n, n}\), the clique in \(L(K_{n, n})\) in the set \(\left\{ \left\{ u, v \right\}: u \in N(v)  \right\} \) (which is a clique) is transitive.   
    \end{enumerate}  
\end{claim}

\begin{explanation}
    Let \(V(K_{n, n}) = X \cupdot Y\), \(X = \left\{ x_1, \dots , x_n \right\} \) and \(Y = \left\{ y_1, \dots , y_n \right\} \). Thus,
    \[
        V(L(K_{n, n})) = E(K_{n, n}) = \left\{ \left\{ x_i, y_j \right\}: 1 \le i, j \le n  \right\},
    \]   
    and 
    \begin{align*}
        E(L(K_{n, n})) &= \left\{ \left\{ \left\{ x_i, y_j \right\}, \left\{ x_{i^{\prime} }, y_j \right\}: 1 \le j \le n, 1 \le i \neq i^{\prime} \le n   \right\} \cup \left\{ \left\{ \left\{ x_i, y_j \right\}, \left\{ x_i, y_{j^{\prime} } \right\}   \right\}: 1 \le i \le n, 1 \le j \neq j^{\prime} \le n  \right\}  \right\} 
    \end{align*}
    Define \(r(x) = x \mod n \in \left\{ 0, 1, \dots , n - 1 \right\} \). Orient \(\left\{ x_i, y_j \right\} - \left\{ x_{i^{\prime} }, y_{j^{\prime} } \right\}  \) as follows:
    \begin{itemize}
        \item If \(i = i^{\prime} \), then \(\left\{ x_i, y_j \right\} \to \left\{ x_i, y_{j^{\prime} } \right\}  \) if \(r(i + j) > r(i + j^{\prime})\). 
        \item If \(j = j^{\prime} \), then \(\left\{ x_i, y_j \right\} \to \left\{ x_{i^{\prime} }, y_j \right\}  \) if \(r(i + j) < r(i^{\prime} + j)\).
    \end{itemize}  
    Then for \(\left\{ x_i, y_j \right\} \), there are \(r(i + j)\) outneighbors of the form \(\left\{ x_i, y_{j^{\prime} } \right\} \) and \((n - 1) - r(i + j)\) outneighbors of the form \(\left\{ x_{i^{\prime} }, y_j \right\} \). Therefore, \(d^+ (\left\{ x_i, y_j \right\} ) = n - 1\), and thus \(\Delta ^+(\vv{D}) = n - 1\). 

    Furthermore, the cliques are transitively indexed in increasing / decreasing order of the remainders: \(\left\{ \left\{ x_i, y_j \right\}: 1 \le j \le n  \right\} \) is transitively ordered from \(j = n - 1 - i \mod n\) to \(j = n - i \mod n\), while \(\left\{ \left\{ x_i, y_j \right\}: 1 \le i \le n  \right\} \) is transitively ordered from \(i = n - j \mod n\) to \(i = n - 1 - j \mod n\).             
\end{explanation}

\begin{claim}
    Any orientation of \(L(K_{n, n})\) in which the cliques 
    \[
        \left\{ \left\{ u, v \right\}: u \in N(v)  \right\}, \quad v \in V(K_{n, n}), 
    \]
    are transitively-oriented, is also kernel-perfect.
\end{claim}

\begin{explanation}
    Given 
    \begin{itemize}
        \item An orientation \(\vv{D}\) of \(L(K_{n, n})\) such that the cliques are transitive. 
        \item \(S \subseteq V(L(K_{n, n}))\) (i.e. \(S \subseteq E(K_{n, n})\)).   
    \end{itemize}
    We want a kernel in \(\vv{D}[S]\). We can use the stable-matching theorem: Let \(K_{n, n} = \left\{ x_1, \dots , x_n \right\} \cup \left\{ y_1, \dots , y_n \right\} \), we build preference lists: for \(x_i\), it prefers \(y_j\) to \(y_{j^{\prime} }\) if 
    \begin{itemize}
        \item \(\left\{ x_i, y_j \right\}, \left\{ x_i, y_{j^{\prime} } \right\} \in S  \) and \(\left\{ x_i, y_j \right\} \leftarrow \left\{ x_i, y_{j^{\prime} } \right\}  \) in \(\vv{D}\), OR 
        \item \(\left\{ x_i, y_j \right\} \in S\) and \(\left\{ x_i, y_{j^{\prime} } \right\} \notin S \), OR 
        \item \(\left\{ x_i, y_j \right\}, \left\{ x_i, y_{j^{\prime}}  \right\} \notin S  \) and \(\left\{ x_i, y_j \right\} \leftarrow \left\{ x_i, y_{j^{\prime} } \right\}   \) in \(\vv{D}\).  
    \end{itemize}
    Note that this is indeed a preference list since all cliques are transitive in the orientation. We can define the preference list for \(y_j\) similarly. Then use Gale-Shapley's algorithm, we can find a stable matching \(M\) of \(K_{n, n}\). Then we show that \(M \cap S\) is a kernel of \(\vv{D} [S]\):
    \begin{enumerate}[(i)]
        \item \(M\) is a matching, so it is an independent set in \(L(K_{n, n})\). 
        \item Suppose \(\left\{ x, y \right\} \in S \setminus M \). Since \(M\) is stable, one of \(x\) or \(y\) prefer their \(M\)-partner; WLOG say \(x\) prefers its \(M\)-partner \(y^{\prime} \) to \(y\). Since \(\left\{ x, y \right\} \in S \), we know \(\left\{ x, y^{\prime}  \right\} \in S \), and so \(\left\{ x, y^{\prime}  \right\} \in S \cap M \) and preferring \(y^{\prime} \). Thus, \(\left\{ x, y \right\} \to \left\{ x, y^{\prime}  \right\}   \) in \(\vv{D}\).                
    \end{enumerate} 
    Now since for all \(S \subseteq L(K_{n, n})\), we can build a kernel in \(\vv{D}[S]\), so \(\vv{D}\) is kernel-perfect.  
\end{explanation}